\documentclass[a4paper,11pt]{article}
\usepackage[T2A]{fontenc}
\usepackage[utf8]{inputenc}
\usepackage[russian]{babel}
\usepackage{amsmath,amssymb,mathtools}
\usepackage{cmbright}
\renewcommand{\familydefault}{\sfdefault}
\usepackage{icomma}
\usepackage[a4paper,left=16mm,right=16mm,top=12mm,bottom=11mm]{geometry}
\usepackage{microtype}
\usepackage{multicol}
\usepackage{xcolor}
\usepackage{hyperref}
\definecolor{accent}{HTML}{15747A}
\definecolor{darkaccent}{HTML}{164C55}
\definecolor{textgray}{HTML}{283238}
\definecolor{muted}{HTML}{65747A}
\definecolor{rulegray}{HTML}{CADADC}
\color{textgray}
\hypersetup{hidelinks}
\setlength{\parindent}{0pt}
\setlength{\parskip}{0pt}
\setlength{\columnsep}{8mm}
\setlength{\abovedisplayskip}{1mm}
\setlength{\belowdisplayskip}{1mm}
\emergencystretch=2em
\newcommand{\sect}[1]{\par\vspace{2.2mm}{\normalsize\bfseries\color{darkaccent}#1}\par\vspace{0.3mm}{\color{rulegray}\rule{\linewidth}{0.45pt}}\par\vspace{0.8mm}}
\newcommand{\fml}[1]{\par\noindent{\normalsize\(\displaystyle #1\)}\par\vspace{2.8mm}}
\begin{document}
\thispagestyle{empty}

\hypersetup{pdftitle={Главные формулы. Володя Хачатурян},pdfauthor={Лёвин Артём Александрович}}
\begin{center}
{\small\bfseries\color{accent}УЧЕБНАЯ ПАМЯТКА \quad\textbullet\quad ОСНОВНЫЕ ФОРМУЛЫ}\par
\vspace{1mm}
{\fontsize{21}{25}\selectfont\bfseries\color{darkaccent} Кинематика и движение по окружности\par}
\vspace{0.8mm}
{\small\color{muted}Володя Хачатурян\quad\textbullet\quad ОГЭ · физика\quad\textbullet\quad 12.10.2026}
\end{center}
\vspace{2mm}
\noindent{\color{accent}\rule{\linewidth}{0.9pt}}
\vspace{1.3mm}
{\bfseries\color{darkaccent}Главные принципы решения}\par
\vspace{1mm}
\begin{multicols}{2}
\raggedcolumns
\small
\noindent\textbullet\quad Проекции скорости и ускорения записывайте с учётом выбранного направления оси.\par\vspace{0.75mm}
\noindent\textbullet\quad Путь находите по всему движению, перемещение — как изменение координаты.\par\vspace{0.75mm}
\columnbreak
\noindent\textbullet\quad Для встречного движения используйте условие равенства координат.\par\vspace{0.75mm}
\noindent\textbullet\quad На одном диске одинакова угловая скорость, у сцепленных колёс — линейная скорость в точке контакта.\par\vspace{0.75mm}
\end{multicols}
\vspace{1.1mm}
\noindent{\color{accent}\rule{\linewidth}{0.5pt}}
\vspace{0.6mm}
\begin{multicols}{2}
\raggedcolumns
\normalsize
\sect{Прямолинейное движение}
\fml{x=x_0+v_{0x}t+\frac{a_xt^2}{2}}
\fml{v_x=v_{0x}+a_xt}
\fml{\Delta x=v_{0x}t+\frac{a_xt^2}{2}}
\sect{Средняя скорость}
\fml{v_{\text{ср}}=\frac{S_{\text{общ}}}{t_{\text{общ}}}}
\fml{v_{\text{сбл}}=v_1+v_2}
\columnbreak
\sect{Окружность}
\fml{\nu=\frac1T,\quad \omega=2\pi\nu}
\fml{v=\omega R=\frac{2\pi R}{T}}
\fml{a_c=\frac{v^2}{R}=\omega^2R}
\sect{Диски и шестерни}
\fml{\omega_1=\omega_2\ \text{(один диск)}}
\fml{v_1=v_2\ \text{(зацепление)}}
\fml{R_1\omega_1=R_2\omega_2}
\end{multicols}
\vfill
{\color{accent}\rule{\linewidth}{0.8pt}}\par\vspace{1.5mm}
{\small\color{muted}Лёвин Артём Александрович эксклюзивно для Володи Хачатуряна}
\end{document}